\section{Esito negativo: archiviazione con disabilitazione del discovery}
\label{sec:esito_negativo}

Quando la strategia non rispetta i criteri di accettazione della issue, non viene integrata nel sistema ma il lavoro svolto \textbf{non viene cancellato}. Viene invece archiviato in modo che:

\begin{itemize}
  \item Il sistema di discovery non carichi la strategia (per evitare di includerla involontariamente in qualche ensemble);
  \item Il codice resti nella repository come riferimento storico;
  \item Le ragioni del fallimento siano documentate, in modo che in futuro non si ri-implementi la stessa strategia partendo da zero.
\end{itemize}

\subsection{Convenzione del prefisso underscore}

La convenzione adottata è la seguente: \textbf{le cartelle e i file il cui nome inizia con il carattere underscore \texttt{\_} vengono ignorati dal sistema di discovery delle strategie.}

In pratica, basta rinominare la cartella della strategia aggiungendo un \texttt{\_} all'inizio del nome per disattivarla dal sistema, lasciando però il contenuto consultabile.

\subsection{Step 13b: Rinomina della cartella della strategia}

Rinominare la cartella aggiungendo l'underscore davanti al nome:

\begin{lstlisting}
git mv strategies/stat_arb_coint strategies/_stat_arb_coint
\end{lstlisting}

Il sistema di discovery, che scansiona \texttt{strategies/} cercando cartelle con \texttt{strategy.py} al loro interno, ignorerà automaticamente quelle che iniziano con \texttt{\_}.

\subsection{Step 14b: Creazione del file \texttt{\_FAILED.md}}

All'interno della cartella rinominata, creare un file \texttt{\_FAILED.md} (anch'esso col prefisso \texttt{\_} per esplicitarne lo stato) che documenta il fallimento. Il file deve contenere:

\begin{itemize}
  \item Razionale originale della strategia (recuperato dalla issue);
  \item Risultati walk-forward OOS effettivamente ottenuti;
  \item Analisi del motivo del fallimento (interpretabile, non solo numerico);
  \item Eventuali iterazioni tentate e perché non hanno portato a un esito positivo;
  \item Raccomandazioni per il futuro (es. "non ri-provare questa stessa strategia su questo universe", oppure "ri-tentare se cambia la disponibilità di dati X").
\end{itemize}

Si raccomanda il seguente template:

\begin{lstlisting}
# Strategia <nome>: archiviata per cattive performance

**Data archiviazione**: YYYY-MM-DD
**Issue di origine**: #<numero>
**Branch di sviluppo**: feature/<nome-strategia>

## Razionale originale
<descrizione dell'edge atteso, con riferimenti a letteratura>

## Criteri di accettazione (dalla issue)
- Sharpe netto OOS > 0.5
- MaxDD < 15%
- Correlazione con momentum < 0.3
- ...

## Risultati walk-forward effettivi
- Sharpe netto OOS: 0.18 (target era >0.5)
- CAGR netto: 2.4% (benchmark SPY: 10.2%)
- Alpha vs SPY: -1.8%
- MaxDD: -22.5%
- Periodo coperto: YYYY-MM-DD -> YYYY-MM-DD

## Analisi del fallimento
<spiegazione interpretabile del perche' la strategia non funziona.
Esempi tipici:
- L'edge in letteratura era documentato su dati pre-2010, post-2010
  e' stato eroso da decimalization e HFT.
- La correlazione con il momentum esistente e' 0.85, quindi non
  porta diversificazione.
- La strategia ha funzionato bene su 2018-2021 ma e' collassata
  nel 2022 (regime di rising rates) suggerendo dipendenza non
  modellata dai tassi.
- Il backtest gross-of-fees era promettente ma i costi di
  transazione hanno mangiato l'edge.>

## Iterazioni tentate
1. <descrizione iterazione 1 e perche' non ha funzionato>
2. <descrizione iterazione 2 e perche' non ha funzionato>
3. ...

## Raccomandazioni
- <indicazione su cosa NON ri-provare>
- <indicazione su cosa potrebbe valere la pena ri-tentare in futuro,
  e a quali condizioni>
\end{lstlisting}

\subsection{Step 15b: Aggiornamento documentazione globale}

Per dare visibilita' a livello di progetto del tentativo fallito, si raccomanda di:

\begin{itemize}
  \item Aggiungere una riga nel registro globale dei tentativi (se presente, ad esempio in \texttt{ROADMAP\_FASE2\_v2.md} o in un eventuale \texttt{JOURNAL.md});
  \item Aggiornare la issue di origine con un commento di chiusura che riassume l'esito e linka al \texttt{\_FAILED.md} appena creato.
\end{itemize}

\subsection{Step 16b: Pull request di archiviazione}

A questo punto si apre una pull request che:

\begin{itemize}
  \item Rinomina la cartella della strategia (commit con \texttt{git mv});
  \item Aggiunge il \texttt{\_FAILED.md} all'interno;
  \item Eventualmente aggiorna documentazione globale.
\end{itemize}

Il template di descrizione della PR di archiviazione e':

\begin{lstlisting}
Closes #<numero-issue>

## Strategia: <nome> [ARCHIVIATA]

## Esito
La strategia non ha superato la validazione walk-forward.
Archiviata in `strategies/_<nome>/` con `_FAILED.md` che
documenta i risultati e le ragioni del fallimento.

## Numeri principali
- Sharpe netto OOS: X.XX (target era >Y.YY)
- CAGR netto: X.XX% (benchmark SPY: Y.YY%)
- MaxDD: -X.XX% (target era < Y.YY%)

## Motivo principale del fallimento
<una riga sintetica>

## Note
Il codice resta nella repository per riferimento storico.
Il sistema di discovery ignora automaticamente la cartella
grazie al prefisso `_`.
\end{lstlisting}

Anche questa PR passa per la normale procedura di review e merge. La differenza rispetto a una PR di esito positivo è semplicemente che il codice mergiato è "archiviato" anziché "attivo".

\subsection{Step 17b: Cleanup}

Identico allo Step 16a: dopo il merge, cancellare il branch remoto e locale.

\subsection{Note operative sul \texttt{\_FAILED.md}}

\begin{enumerate}
  \item Il file deve essere \textbf{onesto}: registrare i numeri effettivi, non versioni "rimaneggiate" per salvare la faccia;
  \item Deve essere \textbf{interpretabile}: un altro sviluppatore che lo legga tra sei mesi deve capire perchè la strategia non ha funzionato;
  \item Deve essere \textbf{actionable}: deve aiutare a non ripetere lo stesso errore. Se il problema era nei dati, indicarlo. Se era nell'edge atteso, indicarlo;
  \item Non e' un \emph{post-mortem} burocratico: l'obiettivo e' costruire memoria collettiva del progetto, non riempire un campo.
\end{enumerate}
